\section{Estado del Arte}
\label{sec:estado_del_arte}

Este capítulo analiza el panorama actual de los modelos de lenguaje, la evolución hacia modelos más eficientes y las investigaciones previas en arquitecturas híbridas.

\subsection{Evolución de los Modelos de Lenguaje: De LLM a SLM}

En los últimos años, el campo del Procesamiento del Lenguaje Natural (NLP) ha estado dominado por los \textit{Large Language Models} (LLM), cuya arquitectura base sigue siendo el Transformer propuesto por Vaswani et al.~\cite{vaswani2017attention}. Modelos como Llama 2~\cite{touvron2023llama} han demostrado capacidades emergentes sorprendentes al escalar el número de parámetros a miles de millones.

Sin embargo, recientemente ha surgido una tendencia hacia los \textit{Small Language Models} (SLM), impulsada por la necesidad de eficiencia y despliegue en entornos con recursos limitados. Investigaciones como \textit{``Textbooks Are All You Need''}~\cite{gunasekar2023textbooks} demuestran que la calidad de los datos puede compensar el tamaño del modelo, permitiendo que arquitecturas de menos de 3 mil millones de parámetros (como Phi-3 o Mistral) compitan en tareas específicas con modelos mucho más grandes. En este contexto, la familia de modelos SmolLM~\cite{allal2024smollm} representa el estado actual de la técnica en SLMs ultra-eficientes, optimizando tanto el proceso de entrenamiento como la curación de datos para maximizar el rendimiento en dispositivos de consumo.

\subsection{Limitaciones de los Transformers Puros}

A pesar de su éxito, los Transformers convencionales presentan un cuello de botella fundamental: la complejidad cuadrática del mecanismo de auto-atención respecto a la longitud de la secuencia~\cite{vaswani2017attention}. Esto implica un consumo de memoria de vídeo (VRAM) y un tiempo de cómputo que crece de forma prohibitiva para contextos largos. En un entorno de producción, esto se traduce en altos costes operativos y latencias elevadas, lo que motiva la búsqueda de arquitecturas con escalado lineal.

\subsection{Fundamentos: Redes Neuronales Recurrentes (RNN)}

Las Redes Neuronales Recurrentes introducidas por Rumelhart et al.~\cite{rumelhart1986learning} representan la primera arquitectura capaz de procesar secuencias manteniendo un estado oculto a lo largo del tiempo. En su forma más simple, una RNN vanilla aplica la misma función de transformación recurrentemente:

\begin{equation}
h_t = \tanh(W_{ih} x_t + W_{hh} h_{t-1} + b)
\end{equation}

donde $h_t$ es el estado oculto en tiempo $t$, $x_t$ es la entrada, y $W_{ih}$, $W_{hh}$ son matrices de pesos. El estado oculto actúa como memoria, permitiendo que la red capture dependencias temporales.

Sin embargo, las RNNs vanilla sufren del problema de \textit{vanishing gradients}: durante el backpropagation a través del tiempo (BPTT), los gradientes se atenúan exponencialmente al propagarse hacia atrás a través de muchos pasos temporales. Esto hace prácticamente imposible entrenar redes que capturen dependencias a largo plazo más allá de 10-20 pasos de tiempo. El cálculo del gradiente requiere multiplicar matrices de Jacobiano repetidamente, y si sus valores propios son menores a 1, el gradiente tiende a cero exponencialmente.

\subsection{LSTM y GRU: Soluciones a Dependencias Largo Plazo}

Para resolver el problema de vanishing gradients, Hochreiter y Schmidhuber propusieron las \textit{Long Short-Term Memory} (LSTM)~\cite{hochreiter1997long} networks. Las LSTM introducen \textit{compuertas} (gates) que controlan qué información fluye a través de la celda, permitiendo que gradientes fluyan sin atenuación:

\begin{align}
f_t &= \sigma(W_f \cdot [h_{t-1}, x_t] + b_f) \quad \text{(forget gate)} \\
i_t &= \sigma(W_i \cdot [h_{t-1}, x_t] + b_i) \quad \text{(input gate)} \\
\tilde{C}_t &= \tanh(W_C \cdot [h_{t-1}, x_t] + b_C) \quad \text{(candidate memory)} \\
C_t &= f_t \odot C_{t-1} + i_t \odot \tilde{C}_t \quad \text{(cell state)} \\
o_t &= \sigma(W_o \cdot [h_{t-1}, x_t] + b_o) \quad \text{(output gate)} \\
h_t &= o_t \odot \tanh(C_t)
\end{align}

Las compuertas utilizan funciones sigmoides que producen valores en $[0,1]$, permitiendo multiplicación suave sin saturación. La célula de memoria $C_t$ mantiene una conexión directa con gradientes a través de suma, evitando desvanecimiento.

Las \textit{Gated Recurrent Units} (GRU) propuestas por Cho et al.~\cite{cho2014learning} simplificaron este mecanismo combinando la compuerta de olvido e input en una sola compuerta de actualización. Las ecuaciones del GRU son:

\begin{align}
r_t &= \sigma(W_r \cdot [h_{t-1}, x_t] + b_r) \quad \text{(reset gate)} \\
z_t &= \sigma(W_z \cdot [h_{t-1}, x_t] + b_z) \quad \text{(update gate)} \\
\tilde{h}_t &= \tanh(W_h \cdot [r_t \odot h_{t-1}, x_t] + b_h) \quad \text{(candidate hidden)} \\
h_t &= (1 - z_t) \odot h_{t-1} + z_t \odot \tilde{h}_t
\end{align}

Las GRUs logran un rendimiento comparable a LSTMs con \textbf{menos parámetros} (~33\% menos) y son más fáciles de entrenar. Ambas arquitecturas eliminan el problema de vanishing gradients al mantener conexiones aditivas directas en la ruta del gradiente, pero siguen siendo secuenciales: el tiempo de inferencia crece linealmente con la longitud de la secuencia.

\subsection{El Mecanismo de Atención y el Nacimiento de los Transformers}

El cuello de botella de la secuencialidad se direccionó mediante el mecanismo de \textit{atención}. Bahdanau et al.~\cite{bahdanau2014neural} propusieron que en lugar de comprimir toda la entrada en un único vector, el decodificador debería poder enfocarse selectivamente en diferentes partes de la entrada en cada paso de decodificación.

El mecanismo de atención calcula un score de relevancia entre una consulta (query) $Q$ y un conjunto de claves (keys) $K$:

\begin{equation}
\text{Attention}(Q, K, V) = \text{softmax}\left(\frac{QK^T}{\sqrt{d_k}}\right) V
\end{equation}

donde $V$ son los valores. La matriz $\text{softmax}(QK^T / \sqrt{d_k})$ actúa como un mapa de pesos que distribuye atención sobre toda la secuencia. A diferencia de la recurrencia, el cálculo de atención es \textbf{paralelizable}: todos los pares de posiciones pueden computarse simultáneamente sin dependencias secuenciales.

Vaswani et al.~\cite{vaswani2017attention} llevaron este concepto al extremo en su modelo \textit{Transformer}, eliminando la recurrencia completamente y construyendo la arquitectura entera sobre mecanismos de auto-atención y redes feed-forward. El Transformer aplica múltiples cabezas de atención (Multi-Head Attention, MHA) en paralelo, permitiendo al modelo atender a diferentes subespacios de representación:

\begin{equation}
\text{MultiHead}(Q, K, V) = \text{Concat}(\text{head}_1, \ldots, \text{head}_h) W^O
\end{equation}

donde cada $\text{head}_i$ es una atención independiente con sus propios pesos. La arquitectura resultante consta de:

\begin{enumerate}
    \item \textbf{Capas de Atención}: Capturan dependencias globales a cualquier distancia en $O(n^2)$ complejidad de tiempo pero paralelizables completamente.
    \item \textbf{Redes Feed-Forward}: Procesan cada posición independientemente con dos transformaciones lineales y una activación ReLU o GELU.
    \item \textbf{Normalización de Capas}: Estabilizan el entrenamiento de redes profundas mediante normalización de activaciones.
    \item \textbf{Conexiones Residuales}: Facilitan el flujo de gradientes a través de capas profundas.
\end{enumerate}

La capacidad del Transformer de procesar en paralelo combinada con su capacidad de atender a cualquier token independientemente de su distancia revolucionó el NLP, permitiendo entrenar en datasets gigantescos y escalar a miles de millones de parámetros. Sin embargo, la complejidad cuadrática $O(n^2)$ del mecanismo de atención respecto a la longitud de secuencia $n$ es prohibitiva para contextos largos.

\subsection{Arquitecturas Híbridas y Recurrencia Moderna}

Para mitigar las limitaciones del Transformer, han resurgido las arquitecturas basadas en recurrencia, pero con enfoques modernos. Las \textit{Gated Recurrent Units} (GRU)~\cite{cho2014learning} ofrecen un manejo eficiente de secuencias con un estado oculto de tamaño fijo, eliminando el coste cuadrático.

Investigaciones recientes como Mamba~\cite{gu2023mamba} han propuesto modelos de espacio de estados (\textit{State Space Models}) que logran un rendimiento comparable a los Transformers pero con inferencia de tiempo lineal. La tendencia actual, y el núcleo de este TFM, es la hibridación: combinar capas de atención para capturar relaciones globales con capas recurrentes o lineales para mantener la eficiencia secuencial.

Esta hibridación se ha visto facilitada por la evolución de los ecosistemas de software, específicamente con la transición de la librería \textit{Transformers} de Hugging Face hacia la versión 5. Mientras que las versiones precedentes (v4.x) requerían implementaciones manuales complejas para integrar componentes recurrentes en el grafo de computación del Transformer, la versión 5 introduce abstracciones nativas para modelos de espacio de estados y arquitecturas secuenciales no atencionales. Estas mejoras permiten una gestión optimizada de la caché de inferencia (KV Cache) y una integración transparente de capas GRU dentro del flujo autorregresivo, garantizando que el modelo mantenga la compatibilidad con herramientas de optimización y despliegue estándar del sector.